\documentstyle[%


righttag%


,11pt%


,amssymb%


,russian%               remove in Eng.


,izvru%                 Set izven% in Eng.


% ,izvru-ks             for brief communications


]{amsart}





%       Math definitions





\newcommand{\mes}{\operatorname{mes}}


\newcommand{\im}{\operatorname{Im}}


\newcommand{\tts}[1]{}





%\hoffset=-15mm%                  For Eng.


 


\setcounter{page}{3}


\god{1999}


\nomer{4~(443)} %                        Remove in Eng.


%\nomer{Vol.\,43, No.\,4}%               Set in Eng.


%\str{\pageref{begin}--\pageref{end}}%  Set in Eng.


\udk{517.544+517.956}





\author{\it Е.Г.\,Ахмедзянова, В.Н.\,Дубинин}


% NB:   ^^ remove in Eng.


\title{Радиальные преобразования множеств \\


и неравенства для трансфинитного диаметра}


\thanks{Данное исследование поддержано Российским фондом фундаментальных


исследований, грант \No\,96-01-00007, а


также --- ISSEP, грант \No\,a97-2120.}





\date{}


%\pagestyle{myheadings}%         Set in Eng.


\pagestyle{plain} %              Remove in Eng.


\begin{document}


%\thispagestyle{plain}%          Set in Eng.


\maketitle


\label{begin}





\section*{Введение}





Пусть $E$ --- ограниченное замкнутое множество в комплексной


плоскости $\bold C_z$.  В 1923 году Фекете ввел понятие трансфинитного


диаметра множества как предел


$$


d(E)=\lim_{n\to\infty}\max\bigg\{


\prod_{1\leq k<l\leq n}|z_k-z_l|^{2/n(n-1)}:


z_1,\dotsc,z_n\in E\bigg\}.


$$


Эта геометрическая характеристика и ее естественные обобщения


имеют многочисленные приложения в теории функций, теории потенциала и


других областях математики (см., напр., [1]--[5]).


Отметим хорошо известные свойства трансфинитного диаметра.


\begin{itemize}


\item[I.] Если $E_1\subset E_2$, то $d(E_1)\leq d(E_2)$.


\item[II.] Пусть функция $w=\phi (z)$ отображает $E$ на $E'$, и 


$\phi (z)$ --- сжимающее


отображение, т.\,е.~$|\phi(z_1)-\phi(z_2)|\leq |z_1-z_2|$ 


для любой пары точек 


$z_1$ и $z_2$ в $E$, тогда $d(E')\leq d(E)$.


\item[III.] Для сегмента $E$ длины $l$ выполняется $d(E)=l/4$.


\item[IV.] Если $p_n(w)=w^n+a_1w^{n-1}+\dotsb+a_n$ --- полином, и


$E'=\{ w: p_n(w)=z,\ z\in E\}$, то $d(E')=(d(E))^{1/n}$.


\item[V.] Пусть $B$ --- связная компонента множества


$\ov {\bold C}_z\setminus E$, содержащая бесконечно


удаленную точку, и $B'$ --- образ $B$ при отображении $w=1/z$.


Тогда $d(E)=(r(B',0))^{-1}$, где $r(B',0)$ --- внутренний радиус области $B'$


относительно точки $z=0$.


\end{itemize}


   


В данной статье рассматриваются задачи, ассоциированные с двумя классическими


нижними оценками трансфинитного диаметра множества $E$. 





Пусть $Q$ --- выпуклое замкнутое множество, и $P_Q E$ ---


проекция множества $E$ на $Q$. Какова нижняя оценка $d(E)$


через линейную меру пересечения $P_Q E\cap\partial Q$\,?


В случае, когда $Q$ --- полуплоскость, ответ вытекает из классического 


результата: трансфинитный диаметр множества $E$  больше либо равен


одной четверти линейной меры ортогональной проекции этого множества


на любую прямую ([1], c.\,294). Отметим, что этот результат проще всего


получить из свойств II и III ([2], c.\,217).


Далее, если $Q$ --- круг, то оценка


следует из классической теоремы Бейрлинга (см., напр.,


[6], c.\,45). Мы изучаем проблему для случая, когда $Q$ представляет 


собой замкнутый угол раствора $\alpha\pi$, $0<\alpha\leq 1$.





Вторая задача восходит к Фекете и касается оценки трансфинитного диаметра


множества $E$ снизу через линейные меры пересечений $E$ с $n$ лучами,


выходящими из заданной точки $z_0$ под равными углами ([7], c.\,117).


Согласно гипотезе Фекете трансфинитный диаметр произвольного 


множества не меньше трансфинитного диаметра множества, состоящего 


из $n$ отрезков,  выходящих из точки $z_0$ под равными углами 


и одинаковой длины, равной среднему геометрическому соответствующих


мер. Доказательство гипотезы Фекете в случае, когда $E$ состоит


из отрезков, выходящих из точки $z_0$, привело Сеге [7]


к понятию ``усредняющей''


симметризации, которая впоследствии была развита Маркусом и нашла


приложения в исследовании других математиков. В общем случае


решение задачи не было известно.


Отметим, что ряд авторов рассматривали задачу Фекете для внутреннего


радиуса области, т.\,е.~в связи со свойством V изучались оценки 


внутреннего радиуса  области $B'$ сверху через линейную меру


пересечения $B'$ с системой из $n$ лучей [8]--[12]. Любопытно сравнить


внутренний радиус ограниченной области и трансфинитный диаметр


ее замыкания. Обе величины при расширении области не уменьшаются,


но при известных радиальных преобразованиях и симметризации [6]


первая величина не уменьшается, а вторая --- не увеличивается.


Косвенно это говорит о сложности проблемы Фекете по сравнению


с соответствующей проблемой для внутреннего радиуса.





Рассмотренные выше оценки трансфинитного диаметра связаны


с радиальным преобразованием замкнутых множеств, при котором эти


множества переходят в множества, звездообразные относительно


некоторой точки так, что трансфинитный диаметр $d(E)$ не увеличивается,


а линейная мера пересечения $E$ с системой радиальных лучей не


уменьшается. В \S\,1 отмечаются частные случаи множеств, когда такое


преобразование возможно, и указывается прием сведения ряда других случаев


к этим частным. Наш подход является реализацией


одного из видов кусочно разделяющей симметризации [6]. Во втором параграфе


мы  обобщаем оценку трансфинитного диаметра множества через линейную меру


проекции этого множества на прямую (теорема 3), решаем задачу Фекете


(неравенство (5)) и даем приложения к мероморфным в круге функциям.





\section{Радиальные преобразования}





Напомним сначала понятие радиального преобразования Маркуса


([8]; [9], c.\,58).


Пусть $E$ --- компакт, и $z_0$ --- произвольная точка плоскости $\bold C_z$.


Для неотрицательного $\rho $ определим 


\begin{gather*}


K_\rho (\theta )=\{ z=z_0+re^{i\theta}: \rho\leq r<\infty \},\quad


0\leq\theta<2\pi, \\


K(\theta )=K_0(\theta );\quad 


l_\rho(\theta, E, z_0)=


\int\limits_{E\cap K_\rho (\theta)}\frac{|dz|}{|z-z_0|},\quad \rho>0.


\end{gather*} 


Положим


$R(\theta )=R(\theta, E, z_0)=\lim\limits_{\rho\to  0}\rho\exp 


l_\rho(\theta,E,z_0)$.


Результатом радиального преобразования множества $E$ относительно


точки $z_0$ называется замкнутое множество 


${\cal R}(E)= \{ z=z_0+re^{i\theta}:0\leq r\leq R(\theta),\ 


0\leq\theta<2\pi\}$,


звездообразное относительно точки $z_0$. Из свойства V и теорем Маркуса


([9], теорема 2.1 или теорема 2.2) вытекает неравенство


\begin{equation}


d(E)\geq d({\cal R}(E)).


\end{equation}





Обозначим через $L(\theta )=L(\theta, E, z_0) $ линейную


меру пересечения $E$ с лучом $K(\theta)$, $0\leq\theta<2\pi$.


Нетрудно показать, что $L(\theta )\geq R(\theta)$ ([8], c.\,625). Пусть


$$


{\cal L}(E)= \{ z=z_0+re^{i\theta}:  0\leq r\leq L(\theta ) ,\ 


0\leq\theta<2\pi\}.


$$


Приведем три частных случая справедливости неравенства 


\begin{equation}


d(E)\geq d({\cal L}(E)).


\end{equation}





\begin{lem} Если множество $E$ расположено на лучах, 


выходящих из некоторой точки $z_0$ под углами, б\'ольшими либо 


равными $\pi/2$, то справедливо неравенство $(2)$.


\end{lem}





\begin{lem} Если множество $E$ расположено на лучах


$K(\theta_1)$, $K(-\theta_1)$, $K(\theta_2)$, $K(-\theta_2)$,


$0\leq\theta_1\leq\theta_2-\pi/2\leq\pi/2$, 


и если $E$ симметрично относительно прямой, проходящей


через точку $z_0$ параллельно действительной оси, то $(2)$ справедливо.


\end{lem}





\begin{lem} Пусть $\theta _1, \theta_2,\dotsc, \theta_n$ ---


произвольные действительные числа и $F$ --- произвольный


компакт на неотрицательной полуоси.Тогда для множества 


$E=\bigcup\limits_{k=1}^{n} \{ z=z_0+re^{i\theta_k}: r\in F\}$


справедливо $(2)$.


\end{lem}





Все три леммы доказываются элементарно с применением свойства


II для сжимающего отображения вида


$\phi (z_0+re^{i\theta})=z_0+\{\mes E\cap[z_0,z_0+re^{i\theta}]\}e^{i\theta}$,


$0\leq r<\infty$, $0\leq\theta\leq 2\pi$.





Далее рассмотрим радиальные преобразования, построенные с применением


линейной меры $L(\theta, E, z_0)$, но отличные от ${\cal L} (E)$.


Для этого понадобится разделяющее преобразование множеств


([6], c.\,29) в одном частном случае. Пусть 


\begin{gather*}


D_k=\{ z: |\arg(z-z_0)-\pi/n-2\pi(k-1)/n|<\pi/n\},\\


\zeta=p_k(z)=z_0+(z-z_0)^{n/2},\ \; z\in D_k,\ \; \im(\zeta-z_0)>0,


\ \; k=1,\dotsc,n.


\end{gather*}


Пусть $E$ --- ограниченное замкнутое множество в плоскости $\bold C_z$, 


удовлетворяющее условию $E\cap\ov D_k\ne\emptyset$, $k=1,\dotsc,n$.


Обозначим через $E_k$ объединение множества $p_k(E\cap\ov D_k)$


c его отражением относительно прямой, проходящей через точку


$z_0$ параллельно действительной оси. Семейство множеств 


$\{ E_k\}^{n}_{k=1}$


называется результатом разделяющего преобразования множества $E$


относительно семейства функций $\{ p_k\} ^{n}_{k=1}$.


Следствие 1.3 ([6], c.\,30) дает


\begin{equation}


d(E)\geq \prod_{k=1}^n(d(E_k))^{2/n^2}.


\end{equation}





Для произвольного множества $E$ и точки $z_0$ рассмотрим звездообразное


множество 


$$


{\cal A}(E)= \{ z=z_0+re^{i\theta}: 0\leq r\leq\sqrt{L(\theta, E,z_0)


L(-\theta, E,z_0)},\ 0\leq\theta<2\pi\}.


$$


Это множество отличается от введенного ранее [11] заменой логарифмической


метрики на линейную.  





\begin{thm} Если множество $E$ расположено на лучах $K(\theta_1)$,


$K(-\theta_1)$, $K(\theta_2)$, $K(-\theta_2)$, 


$0\leq\theta_1\leq\theta_2-\pi/2\leq\pi/2$, то справедливо неравенство


$$


d(E)\geq d({\cal A}(E)).


$$


\end{thm}





\begin{pf} Пусть $\{ E_k\}^{2}_{k=1}$ --- результат разделяющего


преобразования множества $E$ относительно пары функций 


$p_1(z)\equiv z$,


$p_2(z)\equiv 2z_0-z$. Множества $E_1$ и $E_2$ удовлетворяют условиям


леммы 2, согласно которой


$$


d(E_k)\geq d({\cal L}(E_k)),\quad k=1,2.


$$


С учетом (3) имеем 


$$


d(E)\geq \sqrt{d({\cal L}(E_1))d({\cal L}(E_2))}.


$$


Подходящее радиально усредняющее преобразование Маркуса ([9], cc.\,59,~61)


(см. также [6], c.\,32)


ставит в соответствие множествам 


${\cal L} (E_1)$ и


$\wt{\cal L} (E_2):=\{ z=2z_0-\zeta : \zeta \in {\cal L} (E_2)\}$ 


множество ${\cal A} (E)$, при этом  


$\sqrt{d({\cal L}(E_1))d(\wt{\cal L}(E_2))}\geq d({\cal A} (E))$.


\end{pf}





\begin{thm} Если множество


$E$ расположено на лучах $K(\theta+2\pi k/n)$,


$k=1,\dotsc,n$, то справедливо неравенство 


$$


d(E)\geq d({\cal M}(E)),


$$


где через ${\cal M}(E)$ обозначена звезда 


$\{z=z_0+re^{i\theta}:


0\leq r\leq\sqrt[\uproot{2}n]{\prod\limits^n_{k=1}


L(\theta+2\pi k/n, E,z_0)}\} $.


\end{thm}





\begin{pf} Можно считать, что $\theta =0$ и $z_0=0$. 


Обозначим через $\{ E_j\} ^{2n}_{j=1}$


результат разделяющего преобразования множества $E$ относительно семейства


функций $\{p_j\}_{j=1}^{2n}$, $\zeta=p_j (z)\equiv(z^{2n})^{1/2}$,


$z\in D_j$, $\im \zeta >0$, $j=1,\dotsc,2n$.


Согласно (3) имеем


$$


d(E)\geq \prod_{j=1}^{2n}(d(E_j))^{1/2n^2}= 


\prod_{k=1}^{n}(d(E_{2k-1}))^{1/n^2}.


$$ 


Пусть $E'_{2k-1}=\{ w: w^n=z,\ z\in E_{2k-1}\} $. Ввиду свойства IV


$$


d(E'_{2k-1})=(d(E_{2k-1}))^{1/n},\qquad k=1,...,n.


$$


Лемма 3, свойство IV и свойство III дают последовательно


\begin{multline*}


d(E'_{2k-1})\geq d({\cal L} (E'_{2k-1}))=


d([0, L^n(2\pi (k-1)/n, E, 0)])^{1/n}=\\


=(L^n(2\pi (k-1)/n, E, 0)/4)^{1/n},\quad k=1,\dotsc,n.


\end{multline*}


Суммируя выписанные соотношения и вычисляя трансфинитный диаметр 


звезды ${\cal M} (E)$ с помощью свойств III и IV,


получим требуемое неравенство.


\end{pf}





В теоремах 1, 2 речь идет о радиальных преобразованиях, отличных от


${\cal L} (E)$, и кроме того, для специальных множеств $E$.


Однако в этих частных случаях они несут б\'ольшую информацию,


чем преобразования Ахаронова, Кирвана, Маркуса. Например,


если множество $E$ состоит из отрезка $[0, 1-i]$ и произвольного


множества положительной линейной меры $\sqrt{2}$ на луче


$K_{\varepsilon}(\pi /4)$, $\varepsilon >0$, то множество 


${\cal A} (E)$ образовано двумя отрезками


длины $\sqrt{2}$ каждый, а радиальное преобразование


Маркуса дает один отрезок $R(E)=[0,1-i]$  такой же длины.


В этом случае неравенство теоремы 1 сильнее неравенства (1).


Вопрос о справедливости неравенства (2) для произвольных


множеств $E$ остается открытым. 





\section{Приложения}


 


Рассмотрим сначала оценки трансфинитного диаметра, вытекающие 


из свойств радиальных преобразований. Пусть $Q$ --- некоторое 


замкнутое множество плоскости ${\bold C}_z$. Проекцией $P_Q(z)$ точки


$z$ из ${\bold C}_z$ на множество $Q$ называют ближайшую


к $z$ точку множества $Q$. Проекцией множества $E$ на


$Q$ называется множество $P_QE=\{ P_Q(z):z\in E\}$.


Хорошо известно, что оператор проектирования на выпуклое


замкнутое множество является сжимающим отображением.


Учитывая это обстоятельство, из свойства II и теоремы


1 получаем следующий результат.





\begin{thm} Пусть $E$ --- ограниченное замкнутое множество


в плоскости ${\bold C}_z$ и $Q$ --- произвольный замкнутый 


угол этой плоскости раствора $\alpha \pi$, $0<\alpha\leq 1$.


Обозначим через $p$  и $l$  линейные меры пересечений


проекции $P_QE$ со сторонами угла $Q$.


Справедливо неравенство    


\begin{equation}


d(E)\geq \sqrt{pl}(4-2\alpha)^{\alpha/2 -1}


(2\alpha)^{-\alpha/2}.


\end{equation}


Знак равенства достигается для множества $E$, лежащего на границе


угла Q и состоящего из двух равных отрезков, выходящих из


вершины этого угла.


\end{thm}





Отметим, что при подходящем выборе угла $Q$ раствора


$\pi$ $(\alpha =1)$ из неравенства (4) имеем классический


результат


$d(E)\geq l/4$, где $l$ --- линейная мера проекции множества


$E$ на произвольную прямую ([1], c.\,294).  





Из теоремы 2 и свойства I непосредственно вытекает решение задачи Фекете:


для любых $\theta$, $n$ и


$z_0$ справедливо неравенство


\begin{equation}


d(E)\geq\sqrt[\uproot{2}n]


{\textstyle\frac{1}{4}\prod\limits^n_{k=1}L(\theta+2\pi k/n, E,z_0)}.


\end{equation}


Знак равенства достигается для множества $E$, состоящего из $n$


равных отрезков, выходящих из точки $z_0$ под углами 


$\theta+2\pi k/n$, $k=1,\dotsc,n$.





Приложения в теории аналитических функций обусловлены,


в частности, известным результатом Хеймана


\begin{equation}


d(E_f)\leq 1.


\end{equation}


Здесь $E_f$ --- дополнение до множества


значений мероморфной в круге $U=\{ z:|z|<1\} $ функции $w=f(z)$


с разложением в окрестности начала $f(z)=1/z+a_0+a_1z+


a_2z^2+\dotsb$ [13]. Приведем два примера таких приложений.


Для открытого множества $B$, содержащего точку


$z_0$, введем обозначение


$$


S(\theta , B, z_0)=\bigg(\frac{1}{\rho}-\int\limits_{B\cap K_\rho(\theta)}


\frac{|dz|}{|z-z_0|^2}\bigg)^{-1},


$$


где $\rho>0$ настолько мало, что круг $\{ z: |z-z_0|\leq


\rho\}\subset B$. Легко видеть, что 


$S(\theta, B, z_0)$


не зависит от $\rho$.


Кроме того, выполняются неравенства


$$


S(\theta,B,z_0)\leq R(\theta,B,z_0)\leq L(\theta,B,z_0),


$$


где $R(\theta, B, z_0)=\rho\exp\Big(\int\limits_{B\cap K_\rho (\theta)}


\frac{|dz|}{|z-z_0|}\Big)$, а


$L(\theta,B,z_0)$ --- линейная мера пересечения множества $B$


с лучом $K(\theta)$.


Действительно, правое неравенство показано Маркусом ([8], c.\,625).


Для доказательства левого неравенства введем обозначение


$T_\rho=[z_0, z_0+1/\rho e^{-i\theta}]\setminus B'$, где


$B'$ --- образ $B\cap K_\rho(\theta)$ при отображении


$w=z_0+1/(z-z_0)$.


Тогда $S(\theta, B, z_0)


=\lim\limits_{\rho\to 0} (L(-\theta, T_\rho, z_0))^{-1}\leq


\lim\limits_{\rho\to 0} (R(-\theta, T_\rho, z_0))^{-1}=


R(\theta, B, z_0)$.





\begin{thm} Если функция $w=f(z)$ мероморфна в круге $U$ 


и $f(0)=0$, $f'(0)=1$, то для любого действительного


числа $\theta $ и натурального $n$ справедливо неравенство


$$


\prod_{k=1}^n S(\theta+2\pi k/n, f(U), 0)


\geq \frac{1}{4}.


$$


Знак равенства достигается для функции $w=z[1+(e^{-i\theta}z)^n]^{-2/n}$,


отображающей конформно и однолистно круг $U$ на плоскость


${\bf C}_w$ с разрезами по лучам


$\{ w:arg w^n=\theta n,\ |w|\geq \sqrt[\uproot{2}n]{1/4}\} $.


\end{thm}





\begin{pf} Функция $g(z)=1/f(z)$ 


мероморфна в круге $U$, причем в окрестности начала имеет место


разложение $g(z)=1/z


+a_0+a_1z+\dotsb$. Положим $B=f(U)$, $\theta_k=\theta + 2\pi k/n$,


$k=1,\dotsc,n$. Неравенства (6) и (5)


последовательно дают


$$


1\geq d^n(E_g)\geq 


\frac{1}{4} \prod_{k=1}^n L(-\theta_ k, E_g,0)=


\frac{1}{4} \prod_{k=1}^n 


\bigg( \frac{1}{\rho}-\int\limits_{B\cap K_\rho(\theta_k)}


\frac{|dz|}{|z-z_0|^2}\bigg)=


\frac{1}{4} \prod_{k=1}^n S^{-1}(\theta_k, B, 0).


$$


Знак равенства проверяется непосредственно.


\end{pf}





Аналогичные соображения с применением леммы 1 вместо неравенства (5)


и симметризации Штейнера [6] приводят к следующему результату.





\begin{thm} Если функция $w=f(z)$ мероморфна в круге $U$ 


и $f(0)=0$, $f'(0)=1$, то для любого действительного


числа $\theta $ справедливо неравенство


$$


[S^{-1}(\theta,f(U),0)+S^{-1}(\theta+\pi,f(U),0)]^2+


[S^{-1}(\theta+\tfrac{\pi}{2},f(U),0)+S^{-1}(\theta+


\tfrac{3\pi}{2},f(U),0)]^2\leq 16.


$$


Знак равенства имеет место для функции


$w=z[1+c(e^{-i\theta}z)^2+(e^{-i\theta}z)^4]^{-1/2}$,


где


$c$ --- произвольная постоянная, удовлетворяющая условию


$|c|\leq 2$.


\end{thm}





Заметим, что ранее Маркусом были доказаны теоремы 4 и 5, но для регулярных 


функций и с заменой 


$S(\theta,B,0)$ на $R(\theta,B,0)$.


Кроме того, в неравенстве теоремы 5 фигурировали средние геометрические 


вместо средних гармонических, как в нашем случае. 


Таким образом, мы усиливаем соответствующие результаты работ [8], [14].


В заключение отметим, что, применяя неравенство И.П.\,Митюка


(см. [6], c.\,37--38)


вместо неравенства (6), можно получить обобщения теорем


4 и 5 с учетом кратности отображений. 





\begin{thebibliography}{99}





\bibitem{1}


Голузин Г. М. {\it Геометрическая теория функций комплексного переменного.} 


-- М.: Наука, 1966. -- 628 c.


\bibitem{2}


Ландкоф Н.С. {\it Основы современной теории потенциала.} --


М.: Наука, 1966. -- 515~c.


\bibitem{3}


Ransford T. {\it Potential theory in the complex plane}. --


Cambridge: Cambridge Univ. Press, 1995. -- 232~c. 


\bibitem{4}


Overholt M., Schober G. {\it Transfinite extent} // Ann. Acad. Sci.


Fenn. Ser.\,A1. -- 1989. -- V.\,14. -- P.\,277--290.


\bibitem{5}


Amoroso F. {\it $f$-transfinite diameter and number theoretic


applications} //


Ann. Inst. Fourier. -- 1993. -- V.43. -- \No\,4. -- P.\,1179--1198.


\bibitem{6}


Дубинин В.Н. {\it Симметризация в геометрической теории функций


комплексного переменного} // УМН. -- 1994. -- Т.\,49. -- \No\,1.


-- С.\,3--76.


\bibitem{7}


Szego G. {\it On a certain kind of symmetrization and its applications} //


Ann. mat. pura ed appl. -- 1955. -- V.\,40. -- P.\,113--119.


\bibitem{8}


Marcus M. {\it Transformations of domains in the plane and applications


in the theory of functions} // Pacif. J. Math. -- 1964. -- V.\,14.


-- \No\,2. -- P.\,613--626.


\bibitem{9} 


Marcus M. {\it Radial averaging of domains, estimates for Dirichlet


integrals and applications} // J. Anal. Math. -- 1974. -- V.\,27.


-- P.\,47--78. 


\bibitem{10}


Klein M. {\it Estimates for the transfinite diameter with applications to


conformal mapping} // Pacif. J. Math. -- 1967. -- V.\,22. --


\No\,2. -- P.\,267--279.


\bibitem{11}


Aharonov D., Kirwan W.E. {\it A method of symmetrization and applications}.


II // Trans. Amer. Math. Soc. -- 1972. -- V.\,169. -- \No\, 7. -- P.\,279--291.


\bibitem{12}


Zedek M. {\it The linear measures of n-podes in conformal maps of the unit


disk} // J. London Math. Soc. -- 1973. -- V.\,6. -- \No\,2. -- P.\,301--306.


\bibitem{13}


Hayman W.K. {\it Some applications of the transfinite diameter to the


theory of functions} // J. Anal. Math. -- 1951. -- V.\,1. --


P.\,155--179.


\bibitem{14}


Marcus M. {\it Some geometric properties of the image of the 


unit disk by conformal maps} // J. London Math.


Soc. -- 1976. -- V.\,13 -- \No\,1. -- P.\,177--182.





\end{thebibliography}





\vskip 2cm





\post{\it Дальневосточный государственный}{\it Поступила}


\post{\it университет}{17.12.1997}





\label{end}


\end{document}





